% !TeX root = ../main.tex
\section{Grammars and language definitions}

\subsection{Alphabets, strings, and languages}

Before we can define what a \emph{grammar} is, we need the raw material grammars operate
on: strings over an alphabet. This is deliberately the most bare-bones layer of the whole
theory --- no notion of ``meaning'' yet, just symbols and sequences.

\begin{definition}
An \textbf{alphabet} $\Sigma$ is a finite, nonempty set of symbols. A \textbf{string}
(or \textbf{word}) over $\Sigma$ is a finite sequence of symbols from $\Sigma$. The
\textbf{empty string} is written $\varepsilon$, and $\Sigma^*$ denotes the set of
\emph{all} strings over $\Sigma$ (including $\varepsilon$).
\end{definition}

For example, if $\Sigma = \{0, 1\}$, then $\varepsilon, 0, 1, 01, 101, 0011 \in \Sigma^*$.
Note $\Sigma^*$ is always infinite (as long as $\Sigma$ is nonempty), even though $\Sigma$
itself is finite --- this is the first hint of why we need finite machinery (grammars,
automata) to describe infinite sets of strings.

\begin{definition}
A \textbf{language} over $\Sigma$ is any subset $L \subseteq \Sigma^*$.
\end{definition}

That's a deceptively simple definition, and it's worth sitting with it: a language is just
a set of strings, full stop. ``The set of syntactically valid Python programs'' is a
language in this sense (over the alphabet of Unicode characters, or ASCII if you like).
So is ``the set of strings with an equal number of \texttt{a}'s and \texttt{b}'s,'' and so
is ``the set of binary strings representing even numbers.'' Programming languages are
languages in the formal sense, plus a semantics on top --- which is exactly the layering
from the previous section.

\subsubsection{Operations on languages}

Since languages are sets of strings, set operations (union, intersection, complement) apply
directly. But there are two operations specific to strings/languages that are worth
naming explicitly, because they show up constantly in grammars, regular expressions, and
automata:

\begin{itemize}
  \item \textbf{Concatenation.} For strings $x, y$, $xy$ is $x$ followed by $y$. For
    languages, $L_1 L_2 = \{ xy \mid x \in L_1, y \in L_2 \}$ --- every string formed by
    gluing a string from $L_1$ to a string from $L_2$.
  \item \textbf{Kleene star.} $L^* = \bigcup_{i \ge 0} L^i$, where $L^0 = \{\varepsilon\}$
    and $L^{i+1} = L \cdot L^i$. In words: zero or more strings from $L$, concatenated in
    any combination. $\Sigma^*$ itself is exactly the Kleene star of $\Sigma$ viewed as a
    language of length-1 strings.
\end{itemize}

\begin{example}
Let $\Sigma = \{a, b\}$, $L_1 = \{a, ab\}$, $L_2 = \{b, \varepsilon\}$. Then
$L_1 L_2 = \{ab, a, abb, ab\} = \{a, ab, abb\}$ (as a set, duplicates collapse). And
$L_1^* $ contains $\varepsilon, a, ab, aa, aab, aba, \dots$ --- every string you can build by
gluing together zero or more copies of $a$ or $ab$.
\end{example}

Union, concatenation, and Kleene star are the three operations regular expressions are
built from (file 04). Keep them in mind.

\subsection{Grammars: generating languages instead of just naming them}

Set-builder notation (``$L = \{a^n b^n \mid n \ge 0\}$'') is fine for tiny examples, but it
doesn't give you an algorithm for anything --- you can't hand it to a machine and ask
``is this string in $L$?'' or ``generate me a string in $L$.'' A \textbf{grammar} is a
finite, mechanical description of a (possibly infinite) language: a finite set of rewriting
rules that, applied repeatedly, generate exactly the strings of the language.

\begin{definition}
A \textbf{grammar} is a 4-tuple $G = (N, \Sigma, P, S)$ where:
\begin{itemize}
  \item $N$ is a finite set of \textbf{nonterminal} symbols (syntactic categories, like
    ``expression'' or ``statement''),
  \item $\Sigma$ is a finite set of \textbf{terminal} symbols (the alphabet of the actual
    strings we're generating), disjoint from $N$,
  \item $P$ is a finite set of \textbf{production rules}, each of the form $\alpha \to \beta$
    with $\alpha \in (N \cup \Sigma)^*$ containing at least one nonterminal, and
    $\beta \in (N \cup \Sigma)^*$,
  \item $S \in N$ is the \textbf{start symbol}.
\end{itemize}
\end{definition}

The general form above (arbitrary $\alpha$ on the left) is the \emph{unrestricted} grammar,
the top of the Chomsky hierarchy we'll get to shortly. Most of what we do in this course
restricts $\alpha$ to be a single nonterminal, giving context-free grammars (file 05); this
section still states things at the general level so you see where the restriction comes
from.

\subsubsection{Derivations}

A grammar generates strings by \textbf{derivation}: start from $S$, and repeatedly replace
some substring matching the left side of a rule with the rule's right side, until only
terminals remain.

\begin{definition}
Given $\alpha \to \beta \in P$, and a string $\gamma \alpha \delta$, we say
$\gamma \alpha \delta \Rightarrow \gamma \beta \delta$ (``$\gamma\alpha\delta$
\textbf{derives} $\gamma\beta\delta$ in one step''). We write $\Rightarrow^*$ for zero or
more derivation steps. Any string reachable from $S$ via $\Rightarrow^*$ (including $S$
itself) is called a \textbf{sentential form}. The \textbf{language generated by} $G$ is
\[
  L(G) = \{ w \in \Sigma^* \mid S \Rightarrow^* w \}.
\]
\end{definition}

A sentential form may still contain nonterminals (it's a snapshot mid-derivation); a string
in $L(G)$ is a sentential form with \emph{no} nonterminals left. When there's a choice of
which nonterminal to expand next, two conventions matter for later parsing work:

\begin{itemize}
  \item a \textbf{leftmost derivation} always expands the leftmost nonterminal in the
    current sentential form,
  \item a \textbf{rightmost derivation} always expands the rightmost one.
\end{itemize}

These give (in general) different sequences of sentential forms for the same string, even
though the final string and the shape of the parse tree are the same. Leftmost derivations
correspond to top-down parsing; rightmost derivations (in reverse) correspond to bottom-up
parsing --- more on that in the next block.

\subsubsection{A worked example: arithmetic expressions}

Let's generate simple arithmetic expressions over identifiers, \texttt{+}, and \texttt{*},
with parentheses, and actually derive a string.

\[
\begin{aligned}
  N &= \{E\}, \qquad \Sigma = \{ \texttt{id}, +, *, (, ) \}, \qquad S = E \\
  P: \quad & E \to E + E \\
           & E \to E * E \\
           & E \to (E) \\
           & E \to \texttt{id}
\end{aligned}
\]

\begin{example}
Leftmost derivation of $\texttt{id} + \texttt{id} * \texttt{id}$:
\[
\begin{array}{lll}
E & \Rightarrow E + E & (\text{rule } E \to E+E) \\
  & \Rightarrow \texttt{id} + E & (\text{rule } E \to \texttt{id, leftmost } E) \\
  & \Rightarrow \texttt{id} + E * E & (\text{rule } E \to E*E) \\
  & \Rightarrow \texttt{id} + \texttt{id} * E & (\text{rule } E \to \texttt{id}) \\
  & \Rightarrow \texttt{id} + \texttt{id} * \texttt{id} & (\text{rule } E \to \texttt{id})
\end{array}
\]
Every intermediate line ($E+E$, $\texttt{id}+E$, $\texttt{id}+E*E$, \dots) is a sentential
form. The last line has no nonterminals left, so it's a string in $L(G)$.
\end{example}

\subsection{BNF and EBNF}

In practice, nobody writes grammars as bare tuples $(N, \Sigma, P, S)$ --- you write the
production rules in \textbf{Backus-Naur Form} (BNF) or its extended version, \textbf{EBNF},
and everything else ($N$, $\Sigma$, $S$) is inferred from context (nonterminals are
angle-bracketed or capitalized; the start symbol is the one on the left of the first rule,
by convention).

BNF is exactly the productions from above, written with \texttt{::=} instead of $\to$ and
\texttt{|} to combine right-hand sides that share a left-hand side:

\begin{lstlisting}
<expr> ::= <expr> "+" <expr>
         | <expr> "*" <expr>
         | "(" <expr> ")"
         | "id"
\end{lstlisting}

EBNF adds regular-expression-style operators (\texttt{?} for optional, \texttt{*} for zero
or more, \texttt{+} for one or more, grouping with parentheses) purely as \emph{notational
sugar} --- every EBNF rule can be mechanically rewritten into an equivalent set of plain BNF
rules, usually by introducing an extra nonterminal. Here's a small config-language grammar
in EBNF, the kind you'd actually reach for in a real parser:

\begin{lstlisting}
config     ::= entry*
entry      ::= KEY '=' value NEWLINE
value      ::= STRING | NUMBER | list
list       ::= '[' (value (',' value)*)? ']'
\end{lstlisting}

The \texttt{entry*} says ``zero or more entries,'' and \texttt{(',' value)*} says ``zero or
more comma-separated repetitions.'' Written in plain BNF this would need two extra helper
nonterminals (say \texttt{entries} and \texttt{value\_list}) just to express the repetition
recursively. EBNF is strictly more convenient to write and read; it generates exactly the
same class of languages as BNF, nothing more.

\subsection{The Chomsky hierarchy: a map of what's ahead}

Not all grammars are equally hard to parse, and it turns out there's a clean four-level
classification, due to Chomsky, based purely on the shape of the allowed production rules.
This is the map for the rest of this course (and the next block): each level corresponds to
a different class of recognizing machine, with different computational cost.

\begin{center}
\begin{tabular}{p{2.6cm}p{5cm}p{4.3cm}}
\hline
\textbf{Type} & \textbf{Restriction on rules} & \textbf{Recognized by} \\
\hline
Type 3 --- Regular & $A \to a$ or $A \to aB$ (right-linear), $A, B \in N$, $a \in \Sigma$ &
  Finite automata (file 03) \\
Type 2 --- Context-free & $A \to \beta$, single nonterminal on the left &
  Pushdown automata (file 05) \\
Type 1 --- Context-sensitive & $\alpha A \beta \to \alpha \gamma \beta$, $|\gamma| \ge 1$ &
  Linear-bounded automata \\
Type 0 --- Recursively enumerable & no restriction &
  Turing machines \\
\hline
\end{tabular}
\end{center}

Each level is a strict subset of the next: regular $\subsetneq$ context-free $\subsetneq$
context-sensitive $\subsetneq$ recursively enumerable. Every regular language is
context-free (just never use more than one nonterminal ``in flight''), but not every
context-free language is regular (we'll \emph{prove} $\{a^n b^n \mid n \ge 0\}$ isn't
regular in file 04, and it's a totally ordinary context-free language). Most of a
programming language's syntax lives at the context-free level; the things that don't ---
``this variable must be declared before use,'' ``this function call must have the right
number of arguments'' --- are context-sensitive \emph{in spirit} but are essentially never
handled with a real context-sensitive grammar in practice. Instead, they are pushed into a
separate semantic-analysis phase after parsing (next block), precisely because
context-sensitive grammars are a nightmare to parse efficiently. This is the same move we
already saw in the DSL block: keep the grammar simple, push hard constraints to a later
phase.

\subsection{Ambiguity}

\begin{definition}
A grammar $G$ is \textbf{ambiguous} if some string $w \in L(G)$ has two distinct
\emph{parse trees} (equivalently, two distinct leftmost derivations).
\end{definition}

Ambiguity is not a small cosmetic issue --- it means the grammar alone doesn't determine the
structure of the program, so a compiler would have no principled way to decide what a
program means. Two classic examples are worth internalizing because they show up
constantly in real language design.

\subsubsection{Expression precedence}

Take a sloppier version of the arithmetic grammar from before, without separating addition
and multiplication into different nonterminals:

\[
E \to E + E \mid E * E \mid \texttt{id}
\]

The string $\texttt{id} + \texttt{id} * \texttt{id}$ has two different leftmost derivations,
corresponding to two different parse trees:

\begin{center}
\begin{tikzpicture}[level distance=1.1cm, sibling distance=1.6cm, every node/.style={font=\small}]
\node {$E$}
  child { node {$E$} child {node {id}} }
  child { node {$+$} }
  child { node {$E$}
    child { node {$E$} child {node {id}} }
    child { node {$*$} }
    child { node {$E$} child {node {id}} }
  };
\end{tikzpicture}
\hspace{1.5cm}
\begin{tikzpicture}[level distance=1.1cm, sibling distance=1.6cm, every node/.style={font=\small}]
\node {$E$}
  child { node {$E$}
    child { node {$E$} child {node {id}} }
    child { node {$+$} }
    child { node {$E$} child {node {id}} }
  }
  child { node {$*$} }
  child { node {$E$} child {node {id}} };
\end{tikzpicture}
\end{center}

The left tree groups as $\texttt{id} + (\texttt{id} * \texttt{id})$ (multiplication binds
tighter, the reading every programmer expects). The right tree groups as
$(\texttt{id} + \texttt{id}) * \texttt{id}$. Same string, two different computed values
if you evaluated it, and the grammar itself gives no reason to prefer one over the other.

\textbf{The fix} is to encode precedence directly in the grammar, using one nonterminal per
precedence level, lowest precedence at the top:

\[
\begin{aligned}
E &\to E + T \mid T \\
T &\to T * F \mid F \\
F &\to (E) \mid \texttt{id}
\end{aligned}
\]

Now multiplication (built from $T$) can only ever appear \emph{nested inside} an addition
(built from $E$), never the other way around, so the grammar itself forces $*$ to bind
tighter. This grammar is unambiguous, and it's exactly the pattern real language grammars
use for arithmetic expressions.

\subsubsection{Dangling else}

The other classic case:

\[
S \to \mathbf{if}\ b\ \mathbf{then}\ S \mid \mathbf{if}\ b\ \mathbf{then}\ S\ \mathbf{else}\ S \mid \dots
\]

The string $\mathbf{if}\ b_1\ \mathbf{then}\ \mathbf{if}\ b_2\ \mathbf{then}\ S_1\ \mathbf{else}\ S_2$
is ambiguous: the $\mathbf{else}$ could bind to the \emph{inner} $\mathbf{if}$ (the usual
convention, ``else matches the nearest unmatched if'') or to the \emph{outer} one. Both are
valid parse trees under this grammar. We'll come back to this exact example in file 05 and
fix it properly by restructuring the grammar to distinguish ``matched'' from ``unmatched''
if-statements as separate nonterminals --- for now, just notice that the informal rule
(``nearest unmatched if'') that most languages use is a \emph{disambiguation convention}
bolted on top of an ambiguous grammar, rather than a property of an unambiguous one, unless
you do the grammar surgery.

\subsection{Exercises}

\begin{itemize}
  \item Let $\Sigma = \{a, b\}$, $L_1 = \{\varepsilon, a\}$, $L_2 = \{b, bb\}$. Write out
    $L_1 L_2$ and the first four ``levels'' of $L_1^*$ (strings obtainable from 0, 1, 2,
    and 3 concatenated copies of elements of $L_1$).
  \item Give a grammar $G$ (4-tuple, or BNF) that generates
    $L = \{a^n b^n \mid n \ge 0\}$. Give a leftmost derivation of $aabb$ using your grammar,
    showing every sentential form.
  \item Convert the EBNF config-language grammar in this section into plain BNF (no
    \texttt{*}, \texttt{?}, or grouping), introducing helper nonterminals as needed.
  \item Where does $L = \{a^n b^n \mid n \ge 0\}$ sit in the Chomsky hierarchy? Is it
    regular? (You'll be able to prove this rigorously after file 04 --- for now, just give
    your best argument for why no DFA with finitely many states could track $n$ for
    arbitrarily large $n$.)
  \item Take the ambiguous expression grammar $E \to E+E \mid E*E \mid \texttt{id}$ from
    this section and show a string with \emph{three} distinct parse trees (hint: you'll
    need at least two operators of the same kind).
  \item The dangling-else grammar above is ambiguous. Sketch (in prose, you don't need the
    full production rules) how you'd introduce two nonterminals --- one for ``statements
    where every if has a matching else'' and one for ``statements that may have a dangling
    if'' --- to remove the ambiguity. We'll do this properly in file 05.
\end{itemize}
